\documentclass[11pt,a4paper]{article}
\usepackage[utf8]{inputenc}
\usepackage[T1]{fontenc}
\usepackage[catalan]{babel}
\usepackage{lmodern}
\usepackage[sfdefault,lining]{FiraSans}
\usepackage[a4paper,top=2.3cm,bottom=2.2cm,left=1.8cm,right=1.8cm,headheight=1.3cm,headsep=0.4cm,footskip=1cm]{geometry}
\usepackage{amsmath,amssymb}
\usepackage{icomma}
\usepackage{graphicx}
\usepackage[table]{xcolor}
\usepackage{tikz}
\usetikzlibrary{calc}
\usepackage{fontawesome5}
\usepackage{ccicons}
\usepackage{fancyhdr}
\usepackage{lastpage}
\usepackage{tasks}
\usepackage{needspace}
\usepackage[most]{tcolorbox}
\usepackage{array}
\usepackage[hidelinks,pdftitle={Solucions de la fitxa de la unitat 3: Programació lineal. 2n de Batxillerat CCSS},pdfsubject={Matemàtiques Aplicades a les Ciències Socials II},pdfkeywords={programació lineal, inequacions, sistemes d'inequacions, regió factible, vèrtexs, funció objectiu, rectes de nivell, PAU}]{hyperref}

\definecolor{navy}{HTML}{1B2748}
\definecolor{accent}{HTML}{E8603A}
\definecolor{teal}{HTML}{1F9E90}
\definecolor{soft}{HTML}{EEF1F8}
\definecolor{muted}{HTML}{5D6685}

\setlength{\parindent}{0pt}
\setlength{\parskip}{3pt}
\settasks{label=\alph*), label-format=\bfseries\color{navy}, label-width=1.3em, item-indent=1.8em, column-sep=0.6em, after-item-skip=5pt, before-skip=3pt, after-skip=4pt}

% ---------- capçalera i peu ----------
\newcommand{\llicencia}{\href{https://creativecommons.org/licenses/by-nc-sa/4.0/deed.ca}{\ccbyncsa}\ \ Llicència Creative Commons Reconeixement-NoComercial-CompartirIgual 4.0 Internacional (CC BY-NC-SA 4.0)}
\pagestyle{fancy}
\fancyhf{}
\fancyhead[L]{\raisebox{-0.35\height}{\includegraphics[height=0.95cm]{logo-gva.png}}}
\fancyhead[C]{\small\color{muted}Solucions de la fitxa · Unitat 3: Programació lineal · 2n Batxillerat CCSS}
\fancyhead[R]{\raisebox{-0.35\height}{\includegraphics[height=0.95cm]{logo-vermellar.png}}}
\fancyfoot[L]{\scriptsize\color{muted}\llicencia}
\fancyfoot[R]{\small\color{muted}\thepage}
\renewcommand{\headrulewidth}{0.4pt}
\renewcommand{\headrule}{\color{muted!40}\hrule height \headrulewidth}
\fancypagestyle{primera}{\fancyhf{}\renewcommand{\headrulewidth}{0pt}
  \fancyfoot[L]{\scriptsize\color{muted}\llicencia}
  \fancyfoot[R]{\small\color{muted}\thepage}}

% ---------- nivells i exercicis ----------
\newcommand{\estrelles}[1]{{\color{accent}\ifcase#1\or\faStar\faStar[regular]\faStar[regular]\or\faStar\faStar\faStar[regular]\or\faStar\faStar\faStar\fi}}
\newcounter{exer}
\newcommand{\ex}[1]{\par\needspace{4\baselineskip}\medskip\refstepcounter{exer}\noindent{\color{navy}\bfseries\theexer.}\ \ }

\newcommand{\seccio}[2]{\par\needspace{10\baselineskip}\bigskip
  \begin{tcolorbox}[colback=navy,colframe=navy,coltext=white,arc=2mm,boxrule=0pt,left=3mm,right=3mm,top=1.2mm,bottom=1.2mm]
  {\large\bfseries #1\quad #2}
  \end{tcolorbox}\vspace{-2pt}}


% quadre de recordatori (només a la fitxa d'exercicis)
\newcommand{\recorda}[1]{\par\smallskip\begin{tcolorbox}[colback=accent!7,colframe=accent!55,arc=2mm,boxrule=0.5pt,left=2.5mm,right=2.5mm,top=1mm,bottom=1mm]\small{\color{accent}\faLightbulb[regular]}\ \textbf{Recorda:}\ #1\end{tcolorbox}}
% etiqueta dels exercicis de la PAU
\newcommand{\pau}[1]{{\color{teal}\small\bfseries\faUniversity\ PAU · #1.}\ }
\newcommand{\res}[1]{\boxed{#1}}\newcommand{\rest}[1]{\fbox{#1}}

% =====================================================================
% Capçalera FentMates de les fitxes: un color per a cada curs.
% Els colors són els de la pàgina inicial (LEVELS a assets/js/materials.js);
% el logo ve de _eines/logo/pdf/ (python3 _eines/logo/genera.py).
%
% Ús, al final del preàmbul de la unitat:
%   % =====================================================================
% Capçalera FentMates de les fitxes: un color per a cada curs.
% Els colors són els de la pàgina inicial (LEVELS a assets/js/materials.js);
% el logo ve de _eines/logo/pdf/ (python3 _eines/logo/genera.py).
%
% Ús, al final del preàmbul de la unitat:
%   % =====================================================================
% Capçalera FentMates de les fitxes: un color per a cada curs.
% Els colors són els de la pàgina inicial (LEVELS a assets/js/materials.js);
% el logo ve de _eines/logo/pdf/ (python3 _eines/logo/genera.py).
%
% Ús, al final del preàmbul de la unitat:
%   \input{../comu/fentmates.tex}
%   \fmcurs{1eso}
% i, al començament del document:
%   \capcalera{Matemàtiques}{1r ESO}{Fitxa d'exercicis · Unitat 2}{Divisibilitat}
% Opcional abans de \input: \def\fmneed{14} (espai mínim abans de cada secció).
% =====================================================================
\makeatletter
\def\fm@logodir{../../logo/pdf}
\providecommand{\fmneed}{10}
\def\fm@cursnom{}

% clau del curs, color i color del text damunt del color
\newcommand{\fm@curs}[3]{\@namedef{fm@color@#1}{#2}\@namedef{fm@text@#1}{#3}}
\fm@curs{1eso}{E6B84A}{navy}
\fm@curs{2eso}{3B82F6}{white}
\fm@curs{3eso}{FF7A50}{white}
\fm@curs{4eso-a}{A78BFA}{white}
\fm@curs{4eso-b}{6D28D9}{white}
\fm@curs{1batx-mi}{9B1C1C}{white}
\fm@curs{1batx-ccss}{EF6B6B}{white}
\fm@curs{2batx-mii}{1E6B47}{white}
\fm@curs{2batx-ccss}{2BB3A3}{white}

\newcommand{\fmcurs}[1]{%
  \@ifundefined{fm@color@#1}{\PackageError{fentmates}{Curs desconegut: #1}{}}{}%
  \edef\fm@tmp{\noexpand\definecolor{curs}{HTML}{\@nameuse{fm@color@#1}}%
               \noexpand\colorlet{curstext}{\@nameuse{fm@text@#1}}}\fm@tmp
  \colorlet{cursfosc}{curs!72!black}%
  \def\fm@logo{\fm@logodir/vertical-#1.pdf}%
  \def\fm@marca{\fm@logodir/marca-#1.pdf}}

% ---------- capçalera de la primera pàgina ----------
% #1 matèria (petit), #2 curs, #3 tipus de document i unitat, #4 títol
\newcommand{\capcalera}[4]{%
  \gdef\fm@cursnom{#2}%
  \par\noindent
  \begin{tcolorbox}[colback=curs,colframe=curs,boxrule=0pt,arc=2.2mm,width=3.5cm,height=2.15cm,
      valign=center,left=3mm,right=2mm,top=1mm,bottom=1mm,nobeforeafter,box align=center]
    \color{curstext}{\footnotesize #1}\par\vspace{1pt}{\fontsize{19}{21}\selectfont\bfseries #2\par}%
  \end{tcolorbox}\hspace{5mm}%
  \begin{minipage}[c]{\dimexpr\linewidth-3.5cm-5mm-2.9cm\relax}
    {\color{cursfosc}\small\bfseries\MakeUppercase{#3}}\par\vspace{2pt}
    {\color{navy}\fontsize{20}{23}\selectfont\bfseries #4\par}
  \end{minipage}\hfill
  \begin{minipage}[c]{2.6cm}\raggedleft\includegraphics[width=2.6cm]{\fm@logo}\end{minipage}\par
  \vspace{2pt}}

% ---------- barres de secció del color del curs ----------
\renewcommand{\seccio}[2]{\par\needspace{\fmneed\baselineskip}\bigskip
  \begin{tcolorbox}[colback=curs,colframe=curs,coltext=curstext,arc=2mm,boxrule=0pt,left=3mm,right=3mm,top=1.2mm,bottom=1.2mm]
  {\large\bfseries #1\quad #2}
  \end{tcolorbox}\vspace{-2pt}}

% ---------- capçalera i peu de les pàgines ----------
% primera pàgina: logos institucionals al peu
\fancypagestyle{primera}{\fancyhf{}\renewcommand{\headrulewidth}{0pt}
  \fancyfoot[L]{\raisebox{-0.32\height}{\includegraphics[height=0.78cm]{logo-gva.png}}\hspace{5mm}%
    \raisebox{-0.32\height}{\includegraphics[height=0.78cm]{logo-vermellar.png}}}
  \fancyfoot[R]{\scriptsize\color{muted}\href{https://creativecommons.org/licenses/by-nc-sa/4.0/deed.ca}{\ccbyncsa}\ \ CC BY-NC-SA 4.0 · fentmates.org\qquad\small\thepage}}
% la resta: el símbol a l'esquerra i el curs a la dreta, amb el títol al centre com abans
\fancyhead[L]{\raisebox{-0.3\height}{\includegraphics[height=0.62cm]{\fm@marca}}}
\fancyhead[R]{\small\bfseries\color{cursfosc}\fm@cursnom}
\renewcommand{\headrule}{\color{curs!45}\hrule height \headrulewidth}
\makeatother

%   \fmcurs{1eso}
% i, al començament del document:
%   \capcalera{Matemàtiques}{1r ESO}{Fitxa d'exercicis · Unitat 2}{Divisibilitat}
% Opcional abans de \input: \def\fmneed{14} (espai mínim abans de cada secció).
% =====================================================================
\makeatletter
\def\fm@logodir{../../logo/pdf}
\providecommand{\fmneed}{10}
\def\fm@cursnom{}

% clau del curs, color i color del text damunt del color
\newcommand{\fm@curs}[3]{\@namedef{fm@color@#1}{#2}\@namedef{fm@text@#1}{#3}}
\fm@curs{1eso}{E6B84A}{navy}
\fm@curs{2eso}{3B82F6}{white}
\fm@curs{3eso}{FF7A50}{white}
\fm@curs{4eso-a}{A78BFA}{white}
\fm@curs{4eso-b}{6D28D9}{white}
\fm@curs{1batx-mi}{9B1C1C}{white}
\fm@curs{1batx-ccss}{EF6B6B}{white}
\fm@curs{2batx-mii}{1E6B47}{white}
\fm@curs{2batx-ccss}{2BB3A3}{white}

\newcommand{\fmcurs}[1]{%
  \@ifundefined{fm@color@#1}{\PackageError{fentmates}{Curs desconegut: #1}{}}{}%
  \edef\fm@tmp{\noexpand\definecolor{curs}{HTML}{\@nameuse{fm@color@#1}}%
               \noexpand\colorlet{curstext}{\@nameuse{fm@text@#1}}}\fm@tmp
  \colorlet{cursfosc}{curs!72!black}%
  \def\fm@logo{\fm@logodir/vertical-#1.pdf}%
  \def\fm@marca{\fm@logodir/marca-#1.pdf}}

% ---------- capçalera de la primera pàgina ----------
% #1 matèria (petit), #2 curs, #3 tipus de document i unitat, #4 títol
\newcommand{\capcalera}[4]{%
  \gdef\fm@cursnom{#2}%
  \par\noindent
  \begin{tcolorbox}[colback=curs,colframe=curs,boxrule=0pt,arc=2.2mm,width=3.5cm,height=2.15cm,
      valign=center,left=3mm,right=2mm,top=1mm,bottom=1mm,nobeforeafter,box align=center]
    \color{curstext}{\footnotesize #1}\par\vspace{1pt}{\fontsize{19}{21}\selectfont\bfseries #2\par}%
  \end{tcolorbox}\hspace{5mm}%
  \begin{minipage}[c]{\dimexpr\linewidth-3.5cm-5mm-2.9cm\relax}
    {\color{cursfosc}\small\bfseries\MakeUppercase{#3}}\par\vspace{2pt}
    {\color{navy}\fontsize{20}{23}\selectfont\bfseries #4\par}
  \end{minipage}\hfill
  \begin{minipage}[c]{2.6cm}\raggedleft\includegraphics[width=2.6cm]{\fm@logo}\end{minipage}\par
  \vspace{2pt}}

% ---------- barres de secció del color del curs ----------
\renewcommand{\seccio}[2]{\par\needspace{\fmneed\baselineskip}\bigskip
  \begin{tcolorbox}[colback=curs,colframe=curs,coltext=curstext,arc=2mm,boxrule=0pt,left=3mm,right=3mm,top=1.2mm,bottom=1.2mm]
  {\large\bfseries #1\quad #2}
  \end{tcolorbox}\vspace{-2pt}}

% ---------- capçalera i peu de les pàgines ----------
% primera pàgina: logos institucionals al peu
\fancypagestyle{primera}{\fancyhf{}\renewcommand{\headrulewidth}{0pt}
  \fancyfoot[L]{\raisebox{-0.32\height}{\includegraphics[height=0.78cm]{logo-gva.png}}\hspace{5mm}%
    \raisebox{-0.32\height}{\includegraphics[height=0.78cm]{logo-vermellar.png}}}
  \fancyfoot[R]{\scriptsize\color{muted}\href{https://creativecommons.org/licenses/by-nc-sa/4.0/deed.ca}{\ccbyncsa}\ \ CC BY-NC-SA 4.0 · fentmates.org\qquad\small\thepage}}
% la resta: el símbol a l'esquerra i el curs a la dreta, amb el títol al centre com abans
\fancyhead[L]{\raisebox{-0.3\height}{\includegraphics[height=0.62cm]{\fm@marca}}}
\fancyhead[R]{\small\bfseries\color{cursfosc}\fm@cursnom}
\renewcommand{\headrule}{\color{curs!45}\hrule height \headrulewidth}
\makeatother

%   \fmcurs{1eso}
% i, al començament del document:
%   \capcalera{Matemàtiques}{1r ESO}{Fitxa d'exercicis · Unitat 2}{Divisibilitat}
% Opcional abans de \input: \def\fmneed{14} (espai mínim abans de cada secció).
% =====================================================================
\makeatletter
\def\fm@logodir{../../logo/pdf}
\providecommand{\fmneed}{10}
\def\fm@cursnom{}

% clau del curs, color i color del text damunt del color
\newcommand{\fm@curs}[3]{\@namedef{fm@color@#1}{#2}\@namedef{fm@text@#1}{#3}}
\fm@curs{1eso}{E6B84A}{navy}
\fm@curs{2eso}{3B82F6}{white}
\fm@curs{3eso}{FF7A50}{white}
\fm@curs{4eso-a}{A78BFA}{white}
\fm@curs{4eso-b}{6D28D9}{white}
\fm@curs{1batx-mi}{9B1C1C}{white}
\fm@curs{1batx-ccss}{EF6B6B}{white}
\fm@curs{2batx-mii}{1E6B47}{white}
\fm@curs{2batx-ccss}{2BB3A3}{white}

\newcommand{\fmcurs}[1]{%
  \@ifundefined{fm@color@#1}{\PackageError{fentmates}{Curs desconegut: #1}{}}{}%
  \edef\fm@tmp{\noexpand\definecolor{curs}{HTML}{\@nameuse{fm@color@#1}}%
               \noexpand\colorlet{curstext}{\@nameuse{fm@text@#1}}}\fm@tmp
  \colorlet{cursfosc}{curs!72!black}%
  \def\fm@logo{\fm@logodir/vertical-#1.pdf}%
  \def\fm@marca{\fm@logodir/marca-#1.pdf}}

% ---------- capçalera de la primera pàgina ----------
% #1 matèria (petit), #2 curs, #3 tipus de document i unitat, #4 títol
\newcommand{\capcalera}[4]{%
  \gdef\fm@cursnom{#2}%
  \par\noindent
  \begin{tcolorbox}[colback=curs,colframe=curs,boxrule=0pt,arc=2.2mm,width=3.5cm,height=2.15cm,
      valign=center,left=3mm,right=2mm,top=1mm,bottom=1mm,nobeforeafter,box align=center]
    \color{curstext}{\footnotesize #1}\par\vspace{1pt}{\fontsize{19}{21}\selectfont\bfseries #2\par}%
  \end{tcolorbox}\hspace{5mm}%
  \begin{minipage}[c]{\dimexpr\linewidth-3.5cm-5mm-2.9cm\relax}
    {\color{cursfosc}\small\bfseries\MakeUppercase{#3}}\par\vspace{2pt}
    {\color{navy}\fontsize{20}{23}\selectfont\bfseries #4\par}
  \end{minipage}\hfill
  \begin{minipage}[c]{2.6cm}\raggedleft\includegraphics[width=2.6cm]{\fm@logo}\end{minipage}\par
  \vspace{2pt}}

% ---------- barres de secció del color del curs ----------
\renewcommand{\seccio}[2]{\par\needspace{\fmneed\baselineskip}\bigskip
  \begin{tcolorbox}[colback=curs,colframe=curs,coltext=curstext,arc=2mm,boxrule=0pt,left=3mm,right=3mm,top=1.2mm,bottom=1.2mm]
  {\large\bfseries #1\quad #2}
  \end{tcolorbox}\vspace{-2pt}}

% ---------- capçalera i peu de les pàgines ----------
% primera pàgina: logos institucionals al peu
\fancypagestyle{primera}{\fancyhf{}\renewcommand{\headrulewidth}{0pt}
  \fancyfoot[L]{\raisebox{-0.32\height}{\includegraphics[height=0.78cm]{logo-gva.png}}\hspace{5mm}%
    \raisebox{-0.32\height}{\includegraphics[height=0.78cm]{logo-vermellar.png}}}
  \fancyfoot[R]{\scriptsize\color{muted}\href{https://creativecommons.org/licenses/by-nc-sa/4.0/deed.ca}{\ccbyncsa}\ \ CC BY-NC-SA 4.0 · fentmates.org\qquad\small\thepage}}
% la resta: el símbol a l'esquerra i el curs a la dreta, amb el títol al centre com abans
\fancyhead[L]{\raisebox{-0.3\height}{\includegraphics[height=0.62cm]{\fm@marca}}}
\fancyhead[R]{\small\bfseries\color{cursfosc}\fm@cursnom}
\renewcommand{\headrule}{\color{curs!45}\hrule height \headrulewidth}
\makeatother

\fmcurs{2batx-ccss}

\begin{document}
\thispagestyle{primera}

\capcalera{Matemàtiques CCSS}{2n Batx}{Solucions de la fitxa · Unitat 3}{Programació lineal}
\vspace{8pt}

\begin{tcolorbox}[colback=soft,colframe=soft,arc=2mm,left=3mm,right=3mm,top=1.5mm,bottom=1.5mm]
\small Intenta resoldre cada exercici abans de mirar-ne la solució. Pots comprovar les regions factibles, els vèrtexs i els òptims amb el resolutor de programació lineal dels apunts interactius de la unitat.
\end{tcolorbox}

% =====================================================
\seccio{1}{Inequacions lineals amb dues incògnites}

\ex{1} \begin{tasks}(2)
  \task $0\leq4$: \rest{sí}
  \task $6-2=4\leq4$: \rest{sí} (és de la recta frontera)
  \task $-3-6=-9\leq4$: \rest{sí}
  \task $12+4=16\not\leq4$: \rest{no}
\end{tasks}

\ex{1} \begin{tasks}(2)
  \task Recta per $(4,\,0)$ i $(0,\,4)$; \ $(0,\,0)$: $0\leq4$ \checkmark. Semipla de l'origen, recta contínua.\par\begin{tikzpicture}[x=0.4750cm,y=0.4750cm,baseline=(current bounding box.center)]\begin{scope}\clip (-2,-2) rectangle (6,6);\foreach \i in {-2,0,...,6} \draw[muted!25,very thin] (\i,-2)--(\i,6);\foreach \i in {-2,0,...,6} \draw[muted!25,very thin] (-2,\i)--(6,\i);\fill[teal!22] (-2.0000,-2.0000)--(6.0000,-2.0000)--(-2.0000,6.0000)--cycle;\end{scope}\draw[-latex,thick] (-2,0)--(6.5,0) node[below left=-1pt,font=\scriptsize]{$x$};\draw[-latex,thick] (0,-2)--(0,6.5) node[below left=-1pt,font=\scriptsize]{$y$};\node[below,font=\tiny,text=muted] at (-2,0) {$-2$};\node[below,font=\tiny,text=muted] at (2,0) {$2$};\node[below,font=\tiny,text=muted] at (4,0) {$4$};\node[below,font=\tiny,text=muted] at (6,0) {$6$};\node[left,font=\tiny,text=muted] at (0,-2) {$-2$};\node[left,font=\tiny,text=muted] at (0,2) {$2$};\node[left,font=\tiny,text=muted] at (0,4) {$4$};\node[left,font=\tiny,text=muted] at (0,6) {$6$};\begin{scope}\clip (-2,-2) rectangle (6,6);\draw[accent,line width=1.1pt] (7.0000,-3.0000)--(-3.0000,7.0000);\end{scope}\fill[teal] (0,0) circle (2.2pt);\end{tikzpicture}
  \task Recta per $(1,\,0)$ i $(0,\,-2)$; \ $(0,\,0)$: $0>2$ \texttimes. Semipla que no conté l'origen, recta discontínua.\par\begin{tikzpicture}[x=0.5429cm,y=0.5429cm,baseline=(current bounding box.center)]\begin{scope}\clip (-2,-4) rectangle (5,4);\foreach \i in {-2,-1,...,5} \draw[muted!25,very thin] (\i,-4)--(\i,4);\foreach \i in {-4,-2,...,4} \draw[muted!25,very thin] (-2,\i)--(5,\i);\fill[teal!22] (-1.0000,-4.0000)--(5.0000,-4.0000)--(5.0000,4.0000)--(3.0000,4.0000)--cycle;\end{scope}\draw[-latex,thick] (-2,0)--(5.25,0) node[below left=-1pt,font=\scriptsize]{$x$};\draw[-latex,thick] (0,-4)--(0,4.5) node[below left=-1pt,font=\scriptsize]{$y$};\node[below,font=\tiny,text=muted] at (-2,0) {$-2$};\node[below,font=\tiny,text=muted] at (-1,0) {$-1$};\node[below,font=\tiny,text=muted] at (1,0) {$1$};\node[below,font=\tiny,text=muted] at (2,0) {$2$};\node[below,font=\tiny,text=muted] at (3,0) {$3$};\node[below,font=\tiny,text=muted] at (4,0) {$4$};\node[below,font=\tiny,text=muted] at (5,0) {$5$};\node[left,font=\tiny,text=muted] at (0,-4) {$-4$};\node[left,font=\tiny,text=muted] at (0,-2) {$-2$};\node[left,font=\tiny,text=muted] at (0,2) {$2$};\node[left,font=\tiny,text=muted] at (0,4) {$4$};\begin{scope}\clip (-2,-4) rectangle (5,4);\draw[accent,line width=1.1pt,dashed] (-1.5000,-5.0000)--(3.5000,5.0000);\end{scope}\fill[red!70!black] (0,0) circle (2.2pt);\end{tikzpicture}
  \task Recta $y=3x$, que passa per l'origen; amb $(1,\,0)$: $0\geq3$ \texttimes. Semipla de damunt de la recta.\par\begin{tikzpicture}[x=0.6333cm,y=0.5750cm,baseline=(current bounding box.center)]\begin{scope}\clip (-3,-3) rectangle (3,5);\foreach \i in {-3,-2,...,3} \draw[muted!25,very thin] (\i,-3)--(\i,5);\foreach \i in {-2,0,...,5} \draw[muted!25,very thin] (-3,\i)--(3,\i);\fill[teal!22] (-3.0000,-3.0000)--(-1.0000,-3.0000)--(1.6667,5.0000)--(-3.0000,5.0000)--cycle;\end{scope}\draw[-latex,thick] (-3,0)--(3.25,0) node[below left=-1pt,font=\scriptsize]{$x$};\draw[-latex,thick] (0,-3)--(0,5.5) node[below left=-1pt,font=\scriptsize]{$y$};\node[below,font=\tiny,text=muted] at (-3,0) {$-3$};\node[below,font=\tiny,text=muted] at (-2,0) {$-2$};\node[below,font=\tiny,text=muted] at (-1,0) {$-1$};\node[below,font=\tiny,text=muted] at (1,0) {$1$};\node[below,font=\tiny,text=muted] at (2,0) {$2$};\node[below,font=\tiny,text=muted] at (3,0) {$3$};\node[left,font=\tiny,text=muted] at (0,-2) {$-2$};\node[left,font=\tiny,text=muted] at (0,2) {$2$};\node[left,font=\tiny,text=muted] at (0,4) {$4$};\begin{scope}\clip (-3,-3) rectangle (3,5);\draw[accent,line width=1.1pt] (-1.3333,-4.0000)--(2.0000,6.0000);\end{scope}\fill[red!70!black] (1,0) circle (2.2pt);\end{tikzpicture}
  \task Recta per $(6,\,0)$ i $(0,\,-2)$; \ $(0,\,0)$: $0<6$ \checkmark. Semipla de l'origen, recta discontínua.\par\begin{tikzpicture}[x=0.3800cm,y=0.3800cm,baseline=(current bounding box.center)]\begin{scope}\clip (-2,-4) rectangle (8,3);\foreach \i in {-2,0,...,8} \draw[muted!25,very thin] (\i,-4)--(\i,3);\foreach \i in {-4,-3,...,3} \draw[muted!25,very thin] (-2,\i)--(8,\i);\fill[teal!22] (8.0000,0.6667)--(8.0000,3.0000)--(-2.0000,3.0000)--(-2.0000,-2.6667)--cycle;\end{scope}\draw[-latex,thick] (-2,0)--(8.5,0) node[below left=-1pt,font=\scriptsize]{$x$};\draw[-latex,thick] (0,-4)--(0,3.25) node[below left=-1pt,font=\scriptsize]{$y$};\node[below,font=\tiny,text=muted] at (-2,0) {$-2$};\node[below,font=\tiny,text=muted] at (2,0) {$2$};\node[below,font=\tiny,text=muted] at (4,0) {$4$};\node[below,font=\tiny,text=muted] at (6,0) {$6$};\node[below,font=\tiny,text=muted] at (8,0) {$8$};\node[left,font=\tiny,text=muted] at (0,-4) {$-4$};\node[left,font=\tiny,text=muted] at (0,-3) {$-3$};\node[left,font=\tiny,text=muted] at (0,-2) {$-2$};\node[left,font=\tiny,text=muted] at (0,-1) {$-1$};\node[left,font=\tiny,text=muted] at (0,1) {$1$};\node[left,font=\tiny,text=muted] at (0,2) {$2$};\node[left,font=\tiny,text=muted] at (0,3) {$3$};\begin{scope}\clip (-2,-4) rectangle (8,3);\draw[accent,line width=1.1pt,dashed] (-3.0000,-3.0000)--(9.0000,1.0000);\end{scope}\fill[teal] (0,0) circle (2.2pt);\end{tikzpicture}
\end{tasks}

\ex{1} \begin{tasks}(2)
  \task Semipla a la \rest{dreta} de la recta vertical $x=-2$ (contínua).\par\begin{tikzpicture}[x=0.4857cm,y=0.4857cm,baseline=(current bounding box.center)]\begin{scope}\clip (-4,-3) rectangle (3,3);\foreach \i in {-4,-3,...,3} \draw[muted!25,very thin] (\i,-3)--(\i,3);\foreach \i in {-3,-2,...,3} \draw[muted!25,very thin] (-4,\i)--(3,\i);\fill[teal!22] (-2.0000,-3.0000)--(3.0000,-3.0000)--(3.0000,3.0000)--(-2.0000,3.0000)--cycle;\end{scope}\draw[-latex,thick] (-4,0)--(3.25,0) node[below left=-1pt,font=\scriptsize]{$x$};\draw[-latex,thick] (0,-3)--(0,3.25) node[below left=-1pt,font=\scriptsize]{$y$};\node[below,font=\tiny,text=muted] at (-4,0) {$-4$};\node[below,font=\tiny,text=muted] at (-3,0) {$-3$};\node[below,font=\tiny,text=muted] at (-2,0) {$-2$};\node[below,font=\tiny,text=muted] at (-1,0) {$-1$};\node[below,font=\tiny,text=muted] at (1,0) {$1$};\node[below,font=\tiny,text=muted] at (2,0) {$2$};\node[below,font=\tiny,text=muted] at (3,0) {$3$};\node[left,font=\tiny,text=muted] at (0,-3) {$-3$};\node[left,font=\tiny,text=muted] at (0,-2) {$-2$};\node[left,font=\tiny,text=muted] at (0,-1) {$-1$};\node[left,font=\tiny,text=muted] at (0,1) {$1$};\node[left,font=\tiny,text=muted] at (0,2) {$2$};\node[left,font=\tiny,text=muted] at (0,3) {$3$};\begin{scope}\clip (-4,-3) rectangle (3,3);\draw[accent,line width=1.1pt] (-2.0000,-4.0000)--(-2.0000,4.0000);\end{scope}\end{tikzpicture}
  \task Semipla \rest{per davall} de la recta horitzontal $y=3$ (discontínua).\par\begin{tikzpicture}[x=0.5667cm,y=0.5667cm,baseline=(current bounding box.center)]\begin{scope}\clip (-3,-2) rectangle (3,5);\foreach \i in {-3,-2,...,3} \draw[muted!25,very thin] (\i,-2)--(\i,5);\foreach \i in {-2,-1,...,5} \draw[muted!25,very thin] (-3,\i)--(3,\i);\fill[teal!22] (-3.0000,-2.0000)--(3.0000,-2.0000)--(3.0000,3.0000)--(-3.0000,3.0000)--cycle;\end{scope}\draw[-latex,thick] (-3,0)--(3.25,0) node[below left=-1pt,font=\scriptsize]{$x$};\draw[-latex,thick] (0,-2)--(0,5.25) node[below left=-1pt,font=\scriptsize]{$y$};\node[below,font=\tiny,text=muted] at (-3,0) {$-3$};\node[below,font=\tiny,text=muted] at (-2,0) {$-2$};\node[below,font=\tiny,text=muted] at (-1,0) {$-1$};\node[below,font=\tiny,text=muted] at (1,0) {$1$};\node[below,font=\tiny,text=muted] at (2,0) {$2$};\node[below,font=\tiny,text=muted] at (3,0) {$3$};\node[left,font=\tiny,text=muted] at (0,-2) {$-2$};\node[left,font=\tiny,text=muted] at (0,-1) {$-1$};\node[left,font=\tiny,text=muted] at (0,1) {$1$};\node[left,font=\tiny,text=muted] at (0,2) {$2$};\node[left,font=\tiny,text=muted] at (0,3) {$3$};\node[left,font=\tiny,text=muted] at (0,4) {$4$};\node[left,font=\tiny,text=muted] at (0,5) {$5$};\begin{scope}\clip (-3,-2) rectangle (3,5);\draw[accent,line width=1.1pt,dashed] (-4.0000,3.0000)--(4.0000,3.0000);\end{scope}\end{tikzpicture}
  \task $x\leq2{,}5$: semipla a l'\rest{esquerra} de $x=2{,}5$ (contínua).\par\begin{tikzpicture}[x=0.5667cm,y=0.5667cm,baseline=(current bounding box.center)]\begin{scope}\clip (-2,-3) rectangle (4,3);\foreach \i in {-2,-1,...,4} \draw[muted!25,very thin] (\i,-3)--(\i,3);\foreach \i in {-3,-2,...,3} \draw[muted!25,very thin] (-2,\i)--(4,\i);\fill[teal!22] (-2.0000,-3.0000)--(2.5000,-3.0000)--(2.5000,3.0000)--(-2.0000,3.0000)--cycle;\end{scope}\draw[-latex,thick] (-2,0)--(4.25,0) node[below left=-1pt,font=\scriptsize]{$x$};\draw[-latex,thick] (0,-3)--(0,3.25) node[below left=-1pt,font=\scriptsize]{$y$};\node[below,font=\tiny,text=muted] at (-2,0) {$-2$};\node[below,font=\tiny,text=muted] at (-1,0) {$-1$};\node[below,font=\tiny,text=muted] at (1,0) {$1$};\node[below,font=\tiny,text=muted] at (2,0) {$2$};\node[below,font=\tiny,text=muted] at (3,0) {$3$};\node[below,font=\tiny,text=muted] at (4,0) {$4$};\node[left,font=\tiny,text=muted] at (0,-3) {$-3$};\node[left,font=\tiny,text=muted] at (0,-2) {$-2$};\node[left,font=\tiny,text=muted] at (0,-1) {$-1$};\node[left,font=\tiny,text=muted] at (0,1) {$1$};\node[left,font=\tiny,text=muted] at (0,2) {$2$};\node[left,font=\tiny,text=muted] at (0,3) {$3$};\begin{scope}\clip (-2,-3) rectangle (4,3);\draw[accent,line width=1.1pt] (2.5000,-4.0000)--(2.5000,4.0000);\end{scope}\end{tikzpicture}
  \task Multipliquem per $-1$ i \textbf{canvia el sentit}: $y\leq-1$. Semipla \rest{per davall} de $y=-1$.\par\begin{tikzpicture}[x=0.5667cm,y=0.5667cm,baseline=(current bounding box.center)]\begin{scope}\clip (-3,-4) rectangle (3,2);\foreach \i in {-3,-2,...,3} \draw[muted!25,very thin] (\i,-4)--(\i,2);\foreach \i in {-4,-3,...,2} \draw[muted!25,very thin] (-3,\i)--(3,\i);\fill[teal!22] (-3.0000,-4.0000)--(3.0000,-4.0000)--(3.0000,-1.0000)--(-3.0000,-1.0000)--cycle;\end{scope}\draw[-latex,thick] (-3,0)--(3.25,0) node[below left=-1pt,font=\scriptsize]{$x$};\draw[-latex,thick] (0,-4)--(0,2.25) node[below left=-1pt,font=\scriptsize]{$y$};\node[below,font=\tiny,text=muted] at (-3,0) {$-3$};\node[below,font=\tiny,text=muted] at (-2,0) {$-2$};\node[below,font=\tiny,text=muted] at (-1,0) {$-1$};\node[below,font=\tiny,text=muted] at (1,0) {$1$};\node[below,font=\tiny,text=muted] at (2,0) {$2$};\node[below,font=\tiny,text=muted] at (3,0) {$3$};\node[left,font=\tiny,text=muted] at (0,-4) {$-4$};\node[left,font=\tiny,text=muted] at (0,-3) {$-3$};\node[left,font=\tiny,text=muted] at (0,-2) {$-2$};\node[left,font=\tiny,text=muted] at (0,-1) {$-1$};\node[left,font=\tiny,text=muted] at (0,1) {$1$};\node[left,font=\tiny,text=muted] at (0,2) {$2$};\begin{scope}\clip (-3,-4) rectangle (3,2);\draw[accent,line width=1.1pt] (-4.0000,-1.0000)--(4.0000,-1.0000);\end{scope}\end{tikzpicture}
\end{tasks}

\ex{2} \begin{tasks}(2)
  \task Recta per $(0,\,2)$ i $(4,\,0)$: $x+2y=4$. L'origen és de l'ombrejat i la recta és contínua: \ $\res{x+2y\leq4}$
  \task Recta per $(0,\,-1)$ i $(2,\,3)$: $y=2x-1$. L'origen és de l'ombrejat ($0>-1$) i la recta és discontínua: \ $\res{y>2x-1}$, és a dir, $2x-y<1$
\end{tasks}

\ex{2} \begin{tasks}(2)
  \task $\res{2x+3y\leq12}$. Recta per $(6,\,0)$ i $(0,\,4)$; semipla de l'origen.\par\begin{tikzpicture}[x=0.3800cm,y=0.3800cm,baseline=(current bounding box.center)]\begin{scope}\clip (-2,-2) rectangle (8,6);\foreach \i in {-2,0,...,8} \draw[muted!25,very thin] (\i,-2)--(\i,6);\foreach \i in {-2,0,...,6} \draw[muted!25,very thin] (-2,\i)--(8,\i);\fill[teal!22] (-2.0000,-2.0000)--(8.0000,-2.0000)--(8.0000,-1.3333)--(-2.0000,5.3333)--cycle;\end{scope}\draw[-latex,thick] (-2,0)--(8.5,0) node[below left=-1pt,font=\scriptsize]{$x$};\draw[-latex,thick] (0,-2)--(0,6.5) node[below left=-1pt,font=\scriptsize]{$y$};\node[below,font=\tiny,text=muted] at (-2,0) {$-2$};\node[below,font=\tiny,text=muted] at (2,0) {$2$};\node[below,font=\tiny,text=muted] at (4,0) {$4$};\node[below,font=\tiny,text=muted] at (6,0) {$6$};\node[below,font=\tiny,text=muted] at (8,0) {$8$};\node[left,font=\tiny,text=muted] at (0,-2) {$-2$};\node[left,font=\tiny,text=muted] at (0,2) {$2$};\node[left,font=\tiny,text=muted] at (0,4) {$4$};\node[left,font=\tiny,text=muted] at (0,6) {$6$};\begin{scope}\clip (-2,-2) rectangle (8,6);\draw[accent,line width=1.1pt] (-3.0000,6.0000)--(9.0000,-2.0000);\end{scope}\fill[teal] (0,0) circle (2.2pt);\end{tikzpicture}
  \task $\res{x-y\geq2}$. Recta per $(2,\,0)$ i $(0,\,-2)$; $(0,\,0)$ no la compleix: semipla de davall.\par\begin{tikzpicture}[x=0.4750cm,y=0.4750cm,baseline=(current bounding box.center)]\begin{scope}\clip (-2,-4) rectangle (6,4);\foreach \i in {-2,0,...,6} \draw[muted!25,very thin] (\i,-4)--(\i,4);\foreach \i in {-4,-2,...,4} \draw[muted!25,very thin] (-2,\i)--(6,\i);\fill[teal!22] (-2.0000,-4.0000)--(6.0000,-4.0000)--(6.0000,4.0000)--cycle;\end{scope}\draw[-latex,thick] (-2,0)--(6.5,0) node[below left=-1pt,font=\scriptsize]{$x$};\draw[-latex,thick] (0,-4)--(0,4.5) node[below left=-1pt,font=\scriptsize]{$y$};\node[below,font=\tiny,text=muted] at (-2,0) {$-2$};\node[below,font=\tiny,text=muted] at (2,0) {$2$};\node[below,font=\tiny,text=muted] at (4,0) {$4$};\node[below,font=\tiny,text=muted] at (6,0) {$6$};\node[left,font=\tiny,text=muted] at (0,-4) {$-4$};\node[left,font=\tiny,text=muted] at (0,-2) {$-2$};\node[left,font=\tiny,text=muted] at (0,2) {$2$};\node[left,font=\tiny,text=muted] at (0,4) {$4$};\begin{scope}\clip (-2,-4) rectangle (6,4);\draw[accent,line width=1.1pt] (-3.0000,-5.0000)--(7.0000,5.0000);\end{scope}\fill[red!70!black] (0,0) circle (2.2pt);\end{tikzpicture}
\end{tasks}


% =====================================================
\seccio{2}{Sistemes d'inequacions. Regió factible i vèrtexs}

\ex{1} \begin{tasks}(2)
  \task \rest{Acotada}. Vèrtexs: $\res{(0,\,0),\ (5,\,0),\ (2,\,3),\ (0,\,1)}$; \ $(2,\,3)$ és el tall de $x+y=5$ i $x-y=-1$.\par\begin{tikzpicture}[x=0.5867cm,y=0.5867cm,baseline=(current bounding box.center)]\begin{scope}\clip (-1,-1) rectangle (6.5,5.5);\foreach \i in {0,2,...,6.5} \draw[muted!25,very thin] (\i,-1)--(\i,5.5);\foreach \i in {-1,0,...,5.5} \draw[muted!25,very thin] (-1,\i)--(6.5,\i);\fill[teal!22] (5.0000,0.0000)--(2.0000,3.0000)--(0.0000,1.0000)--(0.0000,0.0000)--cycle;\end{scope}\draw[-latex,thick] (-1,0)--(7.0,0) node[below left=-1pt,font=\scriptsize]{$x$};\draw[-latex,thick] (0,-1)--(0,5.75) node[below left=-1pt,font=\scriptsize]{$y$};\node[below,font=\tiny,text=muted] at (2,0) {$2$};\node[below,font=\tiny,text=muted] at (4,0) {$4$};\node[below,font=\tiny,text=muted] at (6,0) {$6$};\node[left,font=\tiny,text=muted] at (0,-1) {$-1$};\node[left,font=\tiny,text=muted] at (0,1) {$1$};\node[left,font=\tiny,text=muted] at (0,2) {$2$};\node[left,font=\tiny,text=muted] at (0,3) {$3$};\node[left,font=\tiny,text=muted] at (0,4) {$4$};\node[left,font=\tiny,text=muted] at (0,5) {$5$};\begin{scope}\clip (-1,-1) rectangle (6.5,5.5);\draw[accent,line width=1.1pt] (7.0000,-2.0000)--(-1.5000,6.5000);\draw[blue!70!black,line width=1.1pt] (-3.0000,-2.0000)--(5.5000,6.5000);\end{scope}\filldraw[fill=white,draw=navy] (5.0000,0.0000) circle (1.5pt);\node[font=\tiny,inner sep=1.5pt,fill=white,fill opacity=.75,text opacity=1] at ($(5.0000,0.0000)+(-17.1:0.36cm)$) {$(5,\,0)$};\filldraw[fill=white,draw=navy] (2.0000,3.0000) circle (1.5pt);\node[font=\tiny,inner sep=1.5pt,fill=white,fill opacity=.75,text opacity=1] at ($(2.0000,3.0000)+(82.9:0.36cm)$) {$(2,\,3)$};\filldraw[fill=white,draw=navy] (0.0000,1.0000) circle (1.5pt);\node[font=\tiny,inner sep=1.5pt,fill=white,fill opacity=.75,text opacity=1] at ($(0.0000,1.0000)+(180.0:0.36cm)$) {$(0,\,1)$};\filldraw[fill=white,draw=navy] (0.0000,0.0000) circle (1.5pt);\node[font=\tiny,inner sep=1.5pt,fill=white,fill opacity=.75,text opacity=1] at ($(0.0000,0.0000)+(-150.3:0.36cm)$) {$(0,\,0)$};\end{tikzpicture}
  \task \rest{No acotada}. Vèrtexs: $\res{(0,\,4),\ (1,\,2),\ (5,\,0)}$; \ $(1,\,2)$ és el tall de $2x+y=4$ i $x+2y=5$.\par\begin{tikzpicture}[x=0.5500cm,y=0.5500cm,baseline=(current bounding box.center)]\begin{scope}\clip (-1,-1) rectangle (7,6);\foreach \i in {0,2,...,7} \draw[muted!25,very thin] (\i,-1)--(\i,6);\foreach \i in {-1,0,...,6} \draw[muted!25,very thin] (-1,\i)--(7,\i);\fill[teal!22] (7.0000,0.0000)--(7.0000,6.0000)--(0.0000,6.0000)--(0.0000,4.0000)--(1.0000,2.0000)--(5.0000,0.0000)--cycle;\end{scope}\draw[-latex,thick] (-1,0)--(7.5,0) node[below left=-1pt,font=\scriptsize]{$x$};\draw[-latex,thick] (0,-1)--(0,6.25) node[below left=-1pt,font=\scriptsize]{$y$};\node[below,font=\tiny,text=muted] at (2,0) {$2$};\node[below,font=\tiny,text=muted] at (4,0) {$4$};\node[below,font=\tiny,text=muted] at (6,0) {$6$};\node[left,font=\tiny,text=muted] at (0,-1) {$-1$};\node[left,font=\tiny,text=muted] at (0,1) {$1$};\node[left,font=\tiny,text=muted] at (0,2) {$2$};\node[left,font=\tiny,text=muted] at (0,3) {$3$};\node[left,font=\tiny,text=muted] at (0,4) {$4$};\node[left,font=\tiny,text=muted] at (0,5) {$5$};\node[left,font=\tiny,text=muted] at (0,6) {$6$};\begin{scope}\clip (-1,-1) rectangle (7,6);\draw[accent,line width=1.1pt] (3.0000,-2.0000)--(-1.5000,7.0000);\draw[blue!70!black,line width=1.1pt] (-2.0000,3.5000)--(8.0000,-1.5000);\end{scope}\filldraw[fill=white,draw=navy] (0.0000,4.0000) circle (1.5pt);\node[font=\tiny,inner sep=1.5pt,fill=white,fill opacity=.75,text opacity=1] at ($(0.0000,4.0000)+(135.0:0.36cm)$) {$(0,\,4)$};\filldraw[fill=white,draw=navy] (1.0000,2.0000) circle (1.5pt);\node[font=\tiny,inner sep=1.5pt,fill=white,fill opacity=.75,text opacity=1] at ($(1.0000,2.0000)+(180.0:0.36cm)$) {$(1,\,2)$};\filldraw[fill=white,draw=navy] (5.0000,0.0000) circle (1.5pt);\node[font=\tiny,inner sep=1.5pt,fill=white,fill opacity=.75,text opacity=1] at ($(5.0000,0.0000)+(-33.7:0.36cm)$) {$(5,\,0)$};\end{tikzpicture}
\end{tasks}

\ex{1} És la franja $1\leq x\leq4$, $0\leq y\leq3$, retallada per la recta $x+y=6$. Vèrtexs: $\res{(1,\,0),\ (4,\,0),\ (4,\,2),\ (3,\,3),\ (1,\,3)}$. El punt $(4,\,3)$ no és vèrtex perquè $4+3>6$.

\begin{center}\begin{tikzpicture}[x=0.7143cm,y=0.7143cm,baseline=(current bounding box.center)]\begin{scope}\clip (-0.5,-1) rectangle (6.5,4.5);\foreach \i in {0,1,...,6.5} \draw[muted!25,very thin] (\i,-1)--(\i,4.5);\foreach \i in {-1,0,...,4.5} \draw[muted!25,very thin] (-0.5,\i)--(6.5,\i);\fill[teal!22] (4.0000,0.0000)--(4.0000,2.0000)--(3.0000,3.0000)--(1.0000,3.0000)--(1.0000,0.0000)--cycle;\end{scope}\draw[-latex,thick] (-0.5,0)--(6.75,0) node[below left=-1pt,font=\scriptsize]{$x$};\draw[-latex,thick] (0,-1)--(0,4.75) node[below left=-1pt,font=\scriptsize]{$y$};\node[below,font=\tiny,text=muted] at (1,0) {$1$};\node[below,font=\tiny,text=muted] at (2,0) {$2$};\node[below,font=\tiny,text=muted] at (3,0) {$3$};\node[below,font=\tiny,text=muted] at (4,0) {$4$};\node[below,font=\tiny,text=muted] at (5,0) {$5$};\node[below,font=\tiny,text=muted] at (6,0) {$6$};\node[left,font=\tiny,text=muted] at (0,-1) {$-1$};\node[left,font=\tiny,text=muted] at (0,1) {$1$};\node[left,font=\tiny,text=muted] at (0,2) {$2$};\node[left,font=\tiny,text=muted] at (0,3) {$3$};\node[left,font=\tiny,text=muted] at (0,4) {$4$};\begin{scope}\clip (-0.5,-1) rectangle (6.5,4.5);\draw[accent,line width=1.1pt] (1.0000,-2.0000)--(1.0000,5.5000);\draw[blue!70!black,line width=1.1pt] (4.0000,-2.0000)--(4.0000,5.5000);\draw[violet,line width=1.1pt] (-1.5000,3.0000)--(7.5000,3.0000);\draw[orange!85!black,line width=1.1pt] (8.0000,-2.0000)--(0.5000,5.5000);\end{scope}\filldraw[fill=white,draw=navy] (4.0000,0.0000) circle (1.5pt);\node[font=\tiny,inner sep=1.5pt,fill=white,fill opacity=.75,text opacity=1] at ($(4.0000,0.0000)+(-48.8:0.36cm)$) {$(4,\,0)$};\filldraw[fill=white,draw=navy] (4.0000,2.0000) circle (1.5pt);\node[font=\tiny,inner sep=1.5pt,fill=white,fill opacity=.75,text opacity=1] at ($(4.0000,2.0000)+(15.9:0.36cm)$) {$(4,\,2)$};\filldraw[fill=white,draw=navy] (3.0000,3.0000) circle (1.5pt);\node[font=\tiny,inner sep=1.5pt,fill=white,fill opacity=.75,text opacity=1] at ($(3.0000,3.0000)+(74.1:0.36cm)$) {$(3,\,3)$};\filldraw[fill=white,draw=navy] (1.0000,3.0000) circle (1.5pt);\node[font=\tiny,inner sep=1.5pt,fill=white,fill opacity=.75,text opacity=1] at ($(1.0000,3.0000)+(138.8:0.36cm)$) {$(1,\,3)$};\filldraw[fill=white,draw=navy] (1.0000,0.0000) circle (1.5pt);\node[font=\tiny,inner sep=1.5pt,fill=white,fill opacity=.75,text opacity=1] at ($(1.0000,0.0000)+(-135.0:0.36cm)$) {$(1,\,0)$};\end{tikzpicture}\end{center}

\ex{2} \begin{tasks}(2)
  \task Les rectes $x+y=2$ i $x+y=1$ són paral·leles i cap punt pot complir les dues inequacions: la regió és \rest{buida}, el sistema \rest{no té solució}.
  \task Franja entre dues rectes paral·leles, tallada per $x=0$ i $x=3$: un paral·lelogram de vèrtexs $\res{(0,\,0),\ (3,\,3),\ (3,\,5),\ (0,\,2)}$.\par\begin{tikzpicture}[x=0.6545cm,y=0.6545cm,baseline=(current bounding box.center)]\begin{scope}\clip (-1,-1) rectangle (4.5,6);\foreach \i in {-1,0,...,4.5} \draw[muted!25,very thin] (\i,-1)--(\i,6);\foreach \i in {-1,0,...,6} \draw[muted!25,very thin] (-1,\i)--(4.5,\i);\fill[teal!22] (0.0000,0.0000)--(3.0000,3.0000)--(3.0000,5.0000)--(0.0000,2.0000)--cycle;\end{scope}\draw[-latex,thick] (-1,0)--(4.75,0) node[below left=-1pt,font=\scriptsize]{$x$};\draw[-latex,thick] (0,-1)--(0,6.25) node[below left=-1pt,font=\scriptsize]{$y$};\node[below,font=\tiny,text=muted] at (-1,0) {$-1$};\node[below,font=\tiny,text=muted] at (1,0) {$1$};\node[below,font=\tiny,text=muted] at (2,0) {$2$};\node[below,font=\tiny,text=muted] at (3,0) {$3$};\node[below,font=\tiny,text=muted] at (4,0) {$4$};\node[left,font=\tiny,text=muted] at (0,-1) {$-1$};\node[left,font=\tiny,text=muted] at (0,1) {$1$};\node[left,font=\tiny,text=muted] at (0,2) {$2$};\node[left,font=\tiny,text=muted] at (0,3) {$3$};\node[left,font=\tiny,text=muted] at (0,4) {$4$};\node[left,font=\tiny,text=muted] at (0,5) {$5$};\node[left,font=\tiny,text=muted] at (0,6) {$6$};\begin{scope}\clip (-1,-1) rectangle (4.5,6);\draw[accent,line width=1.1pt] (-2.0000,-2.0000)--(7.0000,7.0000);\draw[blue!70!black,line width=1.1pt] (-4.0000,-2.0000)--(5.0000,7.0000);\draw[violet,line width=1.1pt] (3.0000,-2.0000)--(3.0000,7.0000);\end{scope}\filldraw[fill=white,draw=navy] (0.0000,0.0000) circle (1.5pt);\node[font=\tiny,inner sep=1.5pt,fill=white,fill opacity=.75,text opacity=1] at ($(0.0000,0.0000)+(-121.0:0.36cm)$) {$(0,\,0)$};\filldraw[fill=white,draw=navy] (3.0000,3.0000) circle (1.5pt);\node[font=\tiny,inner sep=1.5pt,fill=white,fill opacity=.75,text opacity=1] at ($(3.0000,3.0000)+(18.4:0.36cm)$) {$(3,\,3)$};\filldraw[fill=white,draw=navy] (3.0000,5.0000) circle (1.5pt);\node[font=\tiny,inner sep=1.5pt,fill=white,fill opacity=.75,text opacity=1] at ($(3.0000,5.0000)+(59.0:0.36cm)$) {$(3,\,5)$};\filldraw[fill=white,draw=navy] (0.0000,2.0000) circle (1.5pt);\node[font=\tiny,inner sep=1.5pt,fill=white,fill opacity=.75,text opacity=1] at ($(0.0000,2.0000)+(-161.6:0.36cm)$) {$(0,\,2)$};\end{tikzpicture}
\end{tasks}

\ex{2} La recta que passa per $(0,\,0)$ i $(2,\,4)$ és $y=2x$; la que passa per $(2,\,4)$ i $(6,\,0)$ és $x+y=6$. Amb punts de l'interior, com $(3,\,1)$, decidim el sentit: \ $\res{y\leq2x,\quad x+y\leq6,\quad y\geq0}$.

\ex{2} \begin{tasks}(1)
  \task $(1,\,2)$: $5\leq8$ i $5\leq9$: \rest{sí}. \ $(2,\,3)$: $8\leq8$ i $9\leq9$: \rest{sí} (és un vèrtex). \ $(3,\,1)$: $5\leq8$, però $10>9$: \rest{no}.
  \task Vèrtexs: $\res{(0,\,0),\ (3,\,0),\ (2,\,3),\ (0,\,4)}$.\par\begin{tikzpicture}[x=0.6154cm,y=0.6154cm,baseline=(current bounding box.center)]\begin{scope}\clip (-1,-1) rectangle (5.5,5.5);\foreach \i in {-1,0,...,5.5} \draw[muted!25,very thin] (\i,-1)--(\i,5.5);\foreach \i in {-1,0,...,5.5} \draw[muted!25,very thin] (-1,\i)--(5.5,\i);\fill[teal!22] (3.0000,0.0000)--(2.0000,3.0000)--(0.0000,4.0000)--(0.0000,0.0000)--cycle;\end{scope}\draw[-latex,thick] (-1,0)--(5.75,0) node[below left=-1pt,font=\scriptsize]{$x$};\draw[-latex,thick] (0,-1)--(0,5.75) node[below left=-1pt,font=\scriptsize]{$y$};\node[below,font=\tiny,text=muted] at (-1,0) {$-1$};\node[below,font=\tiny,text=muted] at (1,0) {$1$};\node[below,font=\tiny,text=muted] at (2,0) {$2$};\node[below,font=\tiny,text=muted] at (3,0) {$3$};\node[below,font=\tiny,text=muted] at (4,0) {$4$};\node[below,font=\tiny,text=muted] at (5,0) {$5$};\node[left,font=\tiny,text=muted] at (0,-1) {$-1$};\node[left,font=\tiny,text=muted] at (0,1) {$1$};\node[left,font=\tiny,text=muted] at (0,2) {$2$};\node[left,font=\tiny,text=muted] at (0,3) {$3$};\node[left,font=\tiny,text=muted] at (0,4) {$4$};\node[left,font=\tiny,text=muted] at (0,5) {$5$};\begin{scope}\clip (-1,-1) rectangle (5.5,5.5);\draw[accent,line width=1.1pt] (-2.0000,5.0000)--(6.5000,0.7500);\draw[blue!70!black,line width=1.1pt] (3.6667,-2.0000)--(0.8333,6.5000);\end{scope}\filldraw[fill=white,draw=navy] (3.0000,0.0000) circle (1.5pt);\node[font=\tiny,inner sep=1.5pt,fill=white,fill opacity=.75,text opacity=1] at ($(3.0000,0.0000)+(-45.0:0.36cm)$) {$(3,\,0)$};\filldraw[fill=white,draw=navy] (2.0000,3.0000) circle (1.5pt);\node[font=\tiny,inner sep=1.5pt,fill=white,fill opacity=.75,text opacity=1] at ($(2.0000,3.0000)+(59.0:0.36cm)$) {$(2,\,3)$};\filldraw[fill=white,draw=navy] (0.0000,4.0000) circle (1.5pt);\node[font=\tiny,inner sep=1.5pt,fill=white,fill opacity=.75,text opacity=1] at ($(0.0000,4.0000)+(119.1:0.36cm)$) {$(0,\,4)$};\filldraw[fill=white,draw=navy] (0.0000,0.0000) circle (1.5pt);\node[font=\tiny,inner sep=1.5pt,fill=white,fill opacity=.75,text opacity=1] at ($(0.0000,0.0000)+(-125.5:0.36cm)$) {$(0,\,0)$};\end{tikzpicture}
\end{tasks}

\ex{2} És \rest{no acotada} (s'estén cap amunt i cap a la dreta). Vèrtexs: $\res{(1,\,4)}$ i $\res{(3,\,2)}$. El punt $(1,\,2)$, tall de $x=1$ i $y=2$, no és vèrtex perquè $1+2<5$.

\begin{center}\begin{tikzpicture}[x=0.6462cm,y=0.6133cm,baseline=(current bounding box.center)]\begin{scope}\clip (-0.5,-0.5) rectangle (6,7);\foreach \i in {0,1,...,6} \draw[muted!25,very thin] (\i,-0.5)--(\i,7);\foreach \i in {0,2,...,7} \draw[muted!25,very thin] (-0.5,\i)--(6,\i);\fill[teal!22] (6.0000,2.0000)--(6.0000,7.0000)--(1.0000,7.0000)--(1.0000,4.0000)--(3.0000,2.0000)--cycle;\end{scope}\draw[-latex,thick] (-0.5,0)--(6.25,0) node[below left=-1pt,font=\scriptsize]{$x$};\draw[-latex,thick] (0,-0.5)--(0,7.5) node[below left=-1pt,font=\scriptsize]{$y$};\node[below,font=\tiny,text=muted] at (1,0) {$1$};\node[below,font=\tiny,text=muted] at (2,0) {$2$};\node[below,font=\tiny,text=muted] at (3,0) {$3$};\node[below,font=\tiny,text=muted] at (4,0) {$4$};\node[below,font=\tiny,text=muted] at (5,0) {$5$};\node[below,font=\tiny,text=muted] at (6,0) {$6$};\node[left,font=\tiny,text=muted] at (0,2) {$2$};\node[left,font=\tiny,text=muted] at (0,4) {$4$};\node[left,font=\tiny,text=muted] at (0,6) {$6$};\begin{scope}\clip (-0.5,-0.5) rectangle (6,7);\draw[accent,line width=1.1pt] (1.0000,-1.5000)--(1.0000,8.0000);\draw[blue!70!black,line width=1.1pt] (-1.5000,2.0000)--(7.0000,2.0000);\draw[violet,line width=1.1pt] (6.5000,-1.5000)--(-3.0000,8.0000);\end{scope}\filldraw[fill=white,draw=navy] (1.0000,4.0000) circle (1.5pt);\node[font=\tiny,inner sep=1.5pt,fill=white,fill opacity=.75,text opacity=1] at ($(1.0000,4.0000)+(136.5:0.36cm)$) {$(1,\,4)$};\filldraw[fill=white,draw=navy] (3.0000,2.0000) circle (1.5pt);\node[font=\tiny,inner sep=1.5pt,fill=white,fill opacity=.75,text opacity=1] at ($(3.0000,2.0000)+(-43.5:0.36cm)$) {$(3,\,2)$};\end{tikzpicture}\end{center}


% =====================================================
\seccio{3}{Optimització d'una funció objectiu sobre una regió}

\ex{1} \begin{center}\small\begin{tabular}{cc}\hline Vèrtex & $f=3x+2y$\\\hline $(0,\,0)$ & $0$ \ $\leftarrow$\\$(5,\,0)$ & $15$ \ $\leftarrow$\\$(2,\,3)$ & $12$\\$(0,\,1)$ & $2$\\\hline\end{tabular}\end{center}Màxim: $\res{15}$ en $(5,\,0)$. \ Mínim: $\res{0}$ en $(0,\,0)$.

\ex{1} \begin{center}\small\begin{tabular}{ccc}\hline Vèrtex & $f=x+2y$ & $g=2x-y$\\\hline $A(0,\,0)$&0&0\\ $B(6,\,0)$&6&12\\ $C(5,\,3)$&11&7\\ $D(2,\,5)$&12&$-1$\\ $E(0,\,4)$&8&$-4$\\\hline\end{tabular}\end{center}$f$: màxim $\res{12}$ en $D$ i mínim $\res{0}$ en $A$. \ $g$: màxim $\res{12}$ en $B$ i mínim $\res{-4}$ en $E$.

\ex{2} \begin{center}\small\begin{tabular}{cc}\hline Vèrtex & $f=4x+3y$\\\hline $(0,\,4)$ & $12$\\$(1,\,2)$ & $10$ \ $\leftarrow$\\$(5,\,0)$ & $20$\\\hline\end{tabular}\end{center}Mínim: $\res{10}$ en $(1,\,2)$. \ La regió s'estén indefinidament cap amunt i cap a la dreta i $f$ hi creix sense límit: \rest{no té màxim}.

\begin{center}\begin{tikzpicture}[x=0.5250cm,y=0.5250cm,baseline=(current bounding box.center)]\begin{scope}\clip (-1,-1) rectangle (7,6);\foreach \i in {0,2,...,7} \draw[muted!25,very thin] (\i,-1)--(\i,6);\foreach \i in {-1,0,...,6} \draw[muted!25,very thin] (-1,\i)--(7,\i);\fill[teal!22] (7.0000,0.0000)--(7.0000,6.0000)--(0.0000,6.0000)--(0.0000,4.0000)--(1.0000,2.0000)--(5.0000,0.0000)--cycle;\end{scope}\draw[-latex,thick] (-1,0)--(7.5,0) node[below left=-1pt,font=\scriptsize]{$x$};\draw[-latex,thick] (0,-1)--(0,6.25) node[below left=-1pt,font=\scriptsize]{$y$};\node[below,font=\tiny,text=muted] at (2,0) {$2$};\node[below,font=\tiny,text=muted] at (4,0) {$4$};\node[below,font=\tiny,text=muted] at (6,0) {$6$};\node[left,font=\tiny,text=muted] at (0,-1) {$-1$};\node[left,font=\tiny,text=muted] at (0,1) {$1$};\node[left,font=\tiny,text=muted] at (0,2) {$2$};\node[left,font=\tiny,text=muted] at (0,3) {$3$};\node[left,font=\tiny,text=muted] at (0,4) {$4$};\node[left,font=\tiny,text=muted] at (0,5) {$5$};\node[left,font=\tiny,text=muted] at (0,6) {$6$};\begin{scope}\clip (-1,-1) rectangle (7,6);\draw[accent,line width=1.1pt] (3.0000,-2.0000)--(-1.5000,7.0000);\draw[blue!70!black,line width=1.1pt] (-2.0000,3.5000)--(8.0000,-1.5000);\draw[navy,thick,densely dashed] (4.0000,-2.0000)--(-2.7500,7.0000);\end{scope}\filldraw[fill=white,draw=navy] (0.0000,4.0000) circle (1.5pt);\node[font=\tiny,inner sep=1.5pt,fill=white,fill opacity=.75,text opacity=1] at ($(0.0000,4.0000)+(135.0:0.36cm)$) {$(0,\,4)$};\fill[accent] (1.0000,2.0000) circle (2.3pt);\node[font=\tiny,inner sep=1.5pt,fill=white,fill opacity=.75,text opacity=1] at ($(1.0000,2.0000)+(180.0:0.36cm)$) {$(1,\,2)$};\filldraw[fill=white,draw=navy] (5.0000,0.0000) circle (1.5pt);\node[font=\tiny,inner sep=1.5pt,fill=white,fill opacity=.75,text opacity=1] at ($(5.0000,0.0000)+(-33.7:0.36cm)$) {$(5,\,0)$};\end{tikzpicture}\end{center}

\ex{2} Les rectes de nivell $4x+3y=k$ són paral·leles a $4x+3y=0$; en augmentar $k$ es desplacen cap amunt i a la dreta. El primer punt de la regió que toquen és $\res{(1,\,2)}$, amb $k=10$ (el mínim). No n'hi ha un últim: no hi ha màxim.


% =====================================================
\seccio{4}{Problemes de programació lineal}

\ex{1} $x$ = taules, $y$ = cadires. \ $\max\ B(x,y)=40x+20y$ \ amb \ $3x+y\leq60$, \ $x+y\leq30$, \ $x,y\geq0$.

\begin{center}\begin{tikzpicture}[x=0.1158cm,y=0.1158cm,baseline=(current bounding box.center)]\begin{scope}\clip (-3,-3) rectangle (35,35);\foreach \i in {0,10,...,35} \draw[muted!25,very thin] (\i,-3)--(\i,35);\foreach \i in {0,10,...,35} \draw[muted!25,very thin] (-3,\i)--(35,\i);\fill[teal!22] (20.0000,0.0000)--(15.0000,15.0000)--(0.0000,30.0000)--(0.0000,0.0000)--cycle;\end{scope}\draw[-latex,thick] (-3,0)--(37.5,0) node[below left=-1pt,font=\scriptsize]{$x$};\draw[-latex,thick] (0,-3)--(0,37.5) node[below left=-1pt,font=\scriptsize]{$y$};\node[below,font=\tiny,text=muted] at (10,0) {$10$};\node[below,font=\tiny,text=muted] at (20,0) {$20$};\node[below,font=\tiny,text=muted] at (30,0) {$30$};\node[left,font=\tiny,text=muted] at (0,10) {$10$};\node[left,font=\tiny,text=muted] at (0,20) {$20$};\node[left,font=\tiny,text=muted] at (0,30) {$30$};\begin{scope}\clip (-3,-3) rectangle (35,35);\draw[accent,line width=1.1pt] (21.3333,-4.0000)--(8.0000,36.0000);\draw[blue!70!black,line width=1.1pt] (34.0000,-4.0000)--(-6.0000,36.0000);\draw[navy,thick,densely dashed] (24.5000,-4.0000)--(4.5000,36.0000);\end{scope}\filldraw[fill=white,draw=navy] (20.0000,0.0000) circle (1.5pt);\node[font=\tiny,inner sep=1.5pt,fill=white,fill opacity=.75,text opacity=1] at ($(20.0000,0.0000)+(-45.0:0.36cm)$) {$(20,\,0)$};\fill[accent] (15.0000,15.0000) circle (2.3pt);\node[font=\tiny,inner sep=1.5pt,fill=white,fill opacity=.75,text opacity=1] at ($(15.0000,15.0000)+(31.0:0.36cm)$) {$(15,\,15)$};\filldraw[fill=white,draw=navy] (0.0000,30.0000) circle (1.5pt);\node[font=\tiny,inner sep=1.5pt,fill=white,fill opacity=.75,text opacity=1] at ($(0.0000,30.0000)+(115.0:0.36cm)$) {$(0,\,30)$};\filldraw[fill=white,draw=navy] (0.0000,0.0000) circle (1.5pt);\node[font=\tiny,inner sep=1.5pt,fill=white,fill opacity=.75,text opacity=1] at ($(0.0000,0.0000)+(-127.9:0.36cm)$) {$(0,\,0)$};\end{tikzpicture}\end{center}\begin{center}\small\begin{tabular}{cc}\hline Vèrtex & $B=40x+20y$\\\hline $(0,\,0)$ & $0$\\$(20,\,0)$ & $800$\\$(15,\,15)$ & $900$ \ $\leftarrow$\\$(0,\,30)$ & $600$\\\hline\end{tabular}\end{center}Cal fabricar \rest{15 taules i 15 cadires}; benefici màxim $\res{900\text{ €}}$.

\ex{2} $x$ = kg de P, $y$ = kg de Q. \ $\min\ C(x,y)=3x+2y$ \ amb \ $2x+y\geq8$, \ $x+3y\geq9$, \ $x,y\geq0$. La regió \textbf{no és acotada}.

\begin{center}\begin{tikzpicture}[x=0.3667cm,y=0.3667cm,baseline=(current bounding box.center)]\begin{scope}\clip (-1,-1) rectangle (11,10);\foreach \i in {0,2,...,11} \draw[muted!25,very thin] (\i,-1)--(\i,10);\foreach \i in {0,2,...,10} \draw[muted!25,very thin] (-1,\i)--(11,\i);\fill[teal!22] (11.0000,0.0000)--(11.0000,10.0000)--(0.0000,10.0000)--(0.0000,8.0000)--(3.0000,2.0000)--(9.0000,0.0000)--cycle;\end{scope}\draw[-latex,thick] (-1,0)--(11.5,0) node[below left=-1pt,font=\scriptsize]{$x$};\draw[-latex,thick] (0,-1)--(0,10.5) node[below left=-1pt,font=\scriptsize]{$y$};\node[below,font=\tiny,text=muted] at (2,0) {$2$};\node[below,font=\tiny,text=muted] at (4,0) {$4$};\node[below,font=\tiny,text=muted] at (6,0) {$6$};\node[below,font=\tiny,text=muted] at (8,0) {$8$};\node[below,font=\tiny,text=muted] at (10,0) {$10$};\node[left,font=\tiny,text=muted] at (0,2) {$2$};\node[left,font=\tiny,text=muted] at (0,4) {$4$};\node[left,font=\tiny,text=muted] at (0,6) {$6$};\node[left,font=\tiny,text=muted] at (0,8) {$8$};\node[left,font=\tiny,text=muted] at (0,10) {$10$};\begin{scope}\clip (-1,-1) rectangle (11,10);\draw[accent,line width=1.1pt] (5.0000,-2.0000)--(-1.5000,11.0000);\draw[blue!70!black,line width=1.1pt] (-2.0000,3.6667)--(12.0000,-1.0000);\draw[navy,thick,densely dashed] (5.6667,-2.0000)--(-3.0000,11.0000);\end{scope}\filldraw[fill=white,draw=navy] (0.0000,8.0000) circle (1.5pt);\node[font=\tiny,inner sep=1.5pt,fill=white,fill opacity=.75,text opacity=1] at ($(0.0000,8.0000)+(130.6:0.36cm)$) {$(0,\,8)$};\fill[accent] (3.0000,2.0000) circle (2.3pt);\node[font=\tiny,inner sep=1.5pt,fill=white,fill opacity=.75,text opacity=1] at ($(3.0000,2.0000)+(-126.9:0.36cm)$) {$(3,\,2)$};\filldraw[fill=white,draw=navy] (9.0000,0.0000) circle (1.5pt);\node[font=\tiny,inner sep=1.5pt,fill=white,fill opacity=.75,text opacity=1] at ($(9.0000,0.0000)+(-33.7:0.36cm)$) {$(9,\,0)$};\end{tikzpicture}\end{center}\begin{center}\small\begin{tabular}{cc}\hline Vèrtex & $C=3x+2y$\\\hline $(0,\,8)$ & $16$\\$(3,\,2)$ & $13$ \ $\leftarrow$\\$(9,\,0)$ & $27$\\\hline\end{tabular}\end{center}Com que els coeficients són positius, allunyar-se de l'origen només encareix el cost: el mínim existeix. Cal comprar \rest{3 kg de P i 2 kg de Q}; cost mínim $\res{13\text{ €}}$.

\ex{2} $x$, $y$ = milers d'euros en A i en B. Cada miler rendeix 40 € en A i 60 € en B: \ $\max\ R(x,y)=40x+60y$ \ amb \ $x+y\leq30$, \ $x\geq10$, \ $y\leq x$, \ $y\geq0$.

\begin{center}\begin{tikzpicture}[x=0.1278cm,y=0.1278cm,baseline=(current bounding box.center)]\begin{scope}\clip (-2,-2) rectangle (34,20);\foreach \i in {0,5,...,34} \draw[muted!25,very thin] (\i,-2)--(\i,20);\foreach \i in {0,5,...,20} \draw[muted!25,very thin] (-2,\i)--(34,\i);\fill[teal!22] (30.0000,0.0000)--(15.0000,15.0000)--(10.0000,10.0000)--(10.0000,0.0000)--cycle;\end{scope}\draw[-latex,thick] (-2,0)--(35.25,0) node[below left=-1pt,font=\scriptsize]{$x$};\draw[-latex,thick] (0,-2)--(0,21.25) node[below left=-1pt,font=\scriptsize]{$y$};\node[below,font=\tiny,text=muted] at (5,0) {$5$};\node[below,font=\tiny,text=muted] at (10,0) {$10$};\node[below,font=\tiny,text=muted] at (15,0) {$15$};\node[below,font=\tiny,text=muted] at (20,0) {$20$};\node[below,font=\tiny,text=muted] at (25,0) {$25$};\node[below,font=\tiny,text=muted] at (30,0) {$30$};\node[left,font=\tiny,text=muted] at (0,5) {$5$};\node[left,font=\tiny,text=muted] at (0,10) {$10$};\node[left,font=\tiny,text=muted] at (0,15) {$15$};\node[left,font=\tiny,text=muted] at (0,20) {$20$};\begin{scope}\clip (-2,-2) rectangle (34,20);\draw[accent,line width=1.1pt] (33.0000,-3.0000)--(9.0000,21.0000);\draw[blue!70!black,line width=1.1pt] (10.0000,-3.0000)--(10.0000,21.0000);\draw[violet,line width=1.1pt] (-3.0000,-3.0000)--(21.0000,21.0000);\draw[navy,thick,densely dashed] (-3.0000,27.0000)--(35.0000,1.6667);\end{scope}\filldraw[fill=white,draw=navy] (30.0000,0.0000) circle (1.5pt);\node[font=\tiny,inner sep=1.5pt,fill=white,fill opacity=.75,text opacity=1] at ($(30.0000,0.0000)+(-24.4:0.36cm)$) {$(30,\,0)$};\fill[accent] (15.0000,15.0000) circle (2.3pt);\node[font=\tiny,inner sep=1.5pt,fill=white,fill opacity=.75,text opacity=1] at ($(15.0000,15.0000)+(98.1:0.36cm)$) {$(15,\,15)$};\filldraw[fill=white,draw=navy] (10.0000,10.0000) circle (1.5pt);\node[font=\tiny,inner sep=1.5pt,fill=white,fill opacity=.75,text opacity=1] at ($(10.0000,10.0000)+(149.0:0.36cm)$) {$(10,\,10)$};\filldraw[fill=white,draw=navy] (10.0000,0.0000) circle (1.5pt);\node[font=\tiny,inner sep=1.5pt,fill=white,fill opacity=.75,text opacity=1] at ($(10.0000,0.0000)+(-135.0:0.36cm)$) {$(10,\,0)$};\end{tikzpicture}\end{center}\begin{center}\small\begin{tabular}{cc}\hline Vèrtex & $R=40x+60y$\\\hline $(10,\,0)$ & $400$\\$(30,\,0)$ & $1200$\\$(15,\,15)$ & $1500$ \ $\leftarrow$\\$(10,\,10)$ & $1000$\\\hline\end{tabular}\end{center}Ha d'invertir \rest{15\,000 € en A i 15\,000 € en B}; rendiment màxim $\res{1500\text{ €}}$.

\ex{3} $x$ = tones de M1 a T1, $y$ = tones de M1 a T2; aleshores M2 envia $20-x$ a T1 i $25-y$ a T2. \ Restriccions: $0\leq x\leq20$, $0\leq y\leq25$, \ $x+y\leq30$ (M1), \ $(20-x)+(25-y)\leq25\Leftrightarrow x+y\geq20$ (M2). \ Cost: $5x+3y+2(20-x)+4(25-y)=3x-y+140$.

\begin{center}\begin{tikzpicture}[x=0.1680cm,y=0.1533cm,baseline=(current bounding box.center)]\begin{scope}\clip (-2,-2) rectangle (23,28);\foreach \i in {0,5,...,23} \draw[muted!25,very thin] (\i,-2)--(\i,28);\foreach \i in {0,5,...,28} \draw[muted!25,very thin] (-2,\i)--(23,\i);\fill[teal!22] (20.0000,0.0000)--(20.0000,10.0000)--(5.0000,25.0000)--(0.0000,25.0000)--(0.0000,20.0000)--cycle;\end{scope}\draw[-latex,thick] (-2,0)--(24.25,0) node[below left=-1pt,font=\scriptsize]{$x$};\draw[-latex,thick] (0,-2)--(0,29.25) node[below left=-1pt,font=\scriptsize]{$y$};\node[below,font=\tiny,text=muted] at (5,0) {$5$};\node[below,font=\tiny,text=muted] at (10,0) {$10$};\node[below,font=\tiny,text=muted] at (15,0) {$15$};\node[below,font=\tiny,text=muted] at (20,0) {$20$};\node[left,font=\tiny,text=muted] at (0,5) {$5$};\node[left,font=\tiny,text=muted] at (0,10) {$10$};\node[left,font=\tiny,text=muted] at (0,15) {$15$};\node[left,font=\tiny,text=muted] at (0,20) {$20$};\node[left,font=\tiny,text=muted] at (0,25) {$25$};\begin{scope}\clip (-2,-2) rectangle (23,28);\draw[accent,line width=1.1pt] (20.0000,-3.0000)--(20.0000,29.0000);\draw[blue!70!black,line width=1.1pt] (-3.0000,25.0000)--(24.0000,25.0000);\draw[violet,line width=1.1pt] (33.0000,-3.0000)--(1.0000,29.0000);\draw[orange!85!black,line width=1.1pt] (23.0000,-3.0000)--(-9.0000,29.0000);\draw[navy,thick,densely dashed] (-9.3333,-3.0000)--(1.3333,29.0000);\end{scope}\filldraw[fill=white,draw=navy] (20.0000,0.0000) circle (1.5pt);\node[font=\tiny,inner sep=1.5pt,fill=white,fill opacity=.75,text opacity=1] at ($(20.0000,0.0000)+(-53.0:0.36cm)$) {$(20,\,0)$};\filldraw[fill=white,draw=navy] (20.0000,10.0000) circle (1.5pt);\node[font=\tiny,inner sep=1.5pt,fill=white,fill opacity=.75,text opacity=1] at ($(20.0000,10.0000)+(-26.5:0.36cm)$) {$(20,\,10)$};\filldraw[fill=white,draw=navy] (5.0000,25.0000) circle (1.5pt);\node[font=\tiny,inner sep=1.5pt,fill=white,fill opacity=.75,text opacity=1] at ($(5.0000,25.0000)+(116.0:0.36cm)$) {$(5,\,25)$};\fill[accent] (0.0000,25.0000) circle (2.3pt);\node[font=\tiny,inner sep=1.5pt,fill=white,fill opacity=.75,text opacity=1] at ($(0.0000,25.0000)+(137.6:0.36cm)$) {$(0,\,25)$};\filldraw[fill=white,draw=navy] (0.0000,20.0000) circle (1.5pt);\node[font=\tiny,inner sep=1.5pt,fill=white,fill opacity=.75,text opacity=1] at ($(0.0000,20.0000)+(157.9:0.36cm)$) {$(0,\,20)$};\end{tikzpicture}\end{center}\begin{center}\small\begin{tabular}{cc}\hline Vèrtex & $3x-y$\\\hline $(0,\,20)$ & $-20$\\$(0,\,25)$ & $-25$ \ $\leftarrow$\\$(5,\,25)$ & $-10$\\$(20,\,10)$ & $50$\\$(20,\,0)$ & $60$\\\hline\end{tabular}\end{center}Mínim de $3x-y$: $-25$ en $(0,\,25)$. M1 envia \rest{0 t a T1 i 25 t a T2} i M2 envia \rest{20 t a T1 i 0 t a T2}. Cost mínim: $-25+140=\res{115\text{ €}}$.

\ex{2} \pau{C. Valenciana, juliol 2025}$x$ = excursions a la muntanya, $y$ = a la platja. \ $\max\ B=700x+500y$ \ amb \ $2x+2y\leq10$, \ $2x+y\leq8$, \ $x,y\geq0$.

\begin{center}\begin{tikzpicture}[x=0.5846cm,y=0.5846cm,baseline=(current bounding box.center)]\begin{scope}\clip (-0.5,-0.5) rectangle (6,6);\foreach \i in {0,1,...,6} \draw[muted!25,very thin] (\i,-0.5)--(\i,6);\foreach \i in {0,1,...,6} \draw[muted!25,very thin] (-0.5,\i)--(6,\i);\fill[teal!22] (4.0000,0.0000)--(3.0000,2.0000)--(0.0000,5.0000)--(0.0000,0.0000)--cycle;\end{scope}\draw[-latex,thick] (-0.5,0)--(6.25,0) node[below left=-1pt,font=\scriptsize]{$x$};\draw[-latex,thick] (0,-0.5)--(0,6.25) node[below left=-1pt,font=\scriptsize]{$y$};\node[below,font=\tiny,text=muted] at (1,0) {$1$};\node[below,font=\tiny,text=muted] at (2,0) {$2$};\node[below,font=\tiny,text=muted] at (3,0) {$3$};\node[below,font=\tiny,text=muted] at (4,0) {$4$};\node[below,font=\tiny,text=muted] at (5,0) {$5$};\node[below,font=\tiny,text=muted] at (6,0) {$6$};\node[left,font=\tiny,text=muted] at (0,1) {$1$};\node[left,font=\tiny,text=muted] at (0,2) {$2$};\node[left,font=\tiny,text=muted] at (0,3) {$3$};\node[left,font=\tiny,text=muted] at (0,4) {$4$};\node[left,font=\tiny,text=muted] at (0,5) {$5$};\node[left,font=\tiny,text=muted] at (0,6) {$6$};\begin{scope}\clip (-0.5,-0.5) rectangle (6,6);\draw[accent,line width=1.1pt] (6.5000,-1.5000)--(-2.0000,7.0000);\draw[blue!70!black,line width=1.1pt] (4.7500,-1.5000)--(0.5000,7.0000);\draw[navy,thick,densely dashed] (5.5000,-1.5000)--(-0.5714,7.0000);\end{scope}\filldraw[fill=white,draw=navy] (4.0000,0.0000) circle (1.5pt);\node[font=\tiny,inner sep=1.5pt,fill=white,fill opacity=.75,text opacity=1] at ($(4.0000,0.0000)+(-37.9:0.36cm)$) {$(4,\,0)$};\fill[accent] (3.0000,2.0000) circle (2.3pt);\node[font=\tiny,inner sep=1.5pt,fill=white,fill opacity=.75,text opacity=1] at ($(3.0000,2.0000)+(11.3:0.36cm)$) {$(3,\,2)$};\filldraw[fill=white,draw=navy] (0.0000,5.0000) circle (1.5pt);\node[font=\tiny,inner sep=1.5pt,fill=white,fill opacity=.75,text opacity=1] at ($(0.0000,5.0000)+(118.3:0.36cm)$) {$(0,\,5)$};\filldraw[fill=white,draw=navy] (0.0000,0.0000) circle (1.5pt);\node[font=\tiny,inner sep=1.5pt,fill=white,fill opacity=.75,text opacity=1] at ($(0.0000,0.0000)+(-135.0:0.36cm)$) {$(0,\,0)$};\end{tikzpicture}\end{center}\begin{center}\small\begin{tabular}{cc}\hline Vèrtex & $B=700x+500y$\\\hline $(0,\,0)$ & $0$\\$(4,\,0)$ & $2800$\\$(3,\,2)$ & $3100$ \ $\leftarrow$\\$(0,\,5)$ & $2500$\\\hline\end{tabular}\end{center}\begin{tasks}(2)
  \task \rest{3 a la muntanya i 2 a la platja}
  \task $\res{3100\text{ €}}$
\end{tasks}

\ex{2} \pau{C. Valenciana, juny 2025}$x$, $y$ = euros en A i en B. \ $\max\ S=0{,}1x+0{,}07y$ \ amb \ $x+y\leq10\,000$, \ $y\geq2000$, \ $x\leq6000$, \ $x\geq y$.

\begin{center}\small\begin{tabular}{cc}\hline Vèrtex & $S=0{,}1x+0{,}07y$\\\hline $(2000,\,2000)$ & $340$\\$(6000,\,2000)$ & $740$\\$(6000,\,4000)$ & $880$ \ $\leftarrow$\\$(5000,\,5000)$ & $850$\\\hline\end{tabular}\end{center}\begin{tasks}(2)
  \task \rest{6000 € en A i 4000 € en B}
  \task $\res{880}$ seguidors
\end{tasks}

\ex{2} \pau{C. Valenciana, juny 2023}$\min\ C=10x+16y$ \ amb \ $4x+2y\geq8$, \ $6x+20y\geq46$, \ $x+12y\geq12$, \ $x,y\geq0$ (regió no acotada).

\begin{center}\begin{tikzpicture}[x=0.3286cm,y=0.3286cm,baseline=(current bounding box.center)]\begin{scope}\clip (-1,-0.5) rectangle (13,5);\foreach \i in {0,2,...,13} \draw[muted!25,very thin] (\i,-0.5)--(\i,5);\foreach \i in {0,1,...,5} \draw[muted!25,very thin] (-1,\i)--(13,\i);\fill[teal!22] (13.0000,0.0000)--(13.0000,5.0000)--(0.0000,5.0000)--(0.0000,4.0000)--(1.0000,2.0000)--(6.0000,0.5000)--(12.0000,0.0000)--cycle;\end{scope}\draw[-latex,thick] (-1,0)--(13.5,0) node[below left=-1pt,font=\scriptsize]{$x$};\draw[-latex,thick] (0,-0.5)--(0,5.25) node[below left=-1pt,font=\scriptsize]{$y$};\node[below,font=\tiny,text=muted] at (2,0) {$2$};\node[below,font=\tiny,text=muted] at (4,0) {$4$};\node[below,font=\tiny,text=muted] at (6,0) {$6$};\node[below,font=\tiny,text=muted] at (8,0) {$8$};\node[below,font=\tiny,text=muted] at (10,0) {$10$};\node[below,font=\tiny,text=muted] at (12,0) {$12$};\node[left,font=\tiny,text=muted] at (0,1) {$1$};\node[left,font=\tiny,text=muted] at (0,2) {$2$};\node[left,font=\tiny,text=muted] at (0,3) {$3$};\node[left,font=\tiny,text=muted] at (0,4) {$4$};\node[left,font=\tiny,text=muted] at (0,5) {$5$};\begin{scope}\clip (-1,-0.5) rectangle (13,5);\draw[accent,line width=1.1pt] (2.7500,-1.5000)--(-1.0000,6.0000);\draw[blue!70!black,line width=1.1pt] (-2.0000,2.9000)--(14.0000,-1.9000);\draw[violet,line width=1.1pt] (-2.0000,1.1667)--(14.0000,-0.1667);\draw[navy,thick,densely dashed] (-2.0000,3.8750)--(14.0000,-6.1250);\end{scope}\filldraw[fill=white,draw=navy] (0.0000,4.0000) circle (1.5pt);\node[font=\tiny,inner sep=1.5pt,fill=white,fill opacity=.75,text opacity=1] at ($(0.0000,4.0000)+(153.4:0.36cm)$) {$(0,\,4)$};\fill[accent] (1.0000,2.0000) circle (2.3pt);\node[font=\tiny,inner sep=1.5pt,fill=white,fill opacity=.75,text opacity=1] at ($(1.0000,2.0000)+(174.3:0.36cm)$) {$(1,\,2)$};\filldraw[fill=white,draw=navy] (6.0000,0.5000) circle (1.5pt);\node[font=\tiny,inner sep=1.5pt,fill=white,fill opacity=.75,text opacity=1] at ($(6.0000,0.5000)+(-42.0:0.36cm)$) {$(6;0{,}5)$};\filldraw[fill=white,draw=navy] (12.0000,0.0000) circle (1.5pt);\node[font=\tiny,inner sep=1.5pt,fill=white,fill opacity=.75,text opacity=1] at ($(12.0000,0.0000)+(-12.6:0.36cm)$) {$(12,\,0)$};\end{tikzpicture}\end{center}\begin{center}\small\begin{tabular}{cc}\hline Vèrtex & $C=10x+16y$\\\hline $(0,\,4)$ & $64$\\$(1,\,2)$ & $42$ \ $\leftarrow$\\$(6;\ 0{,}5)$ & $68$\\$(12,\,0)$ & $120$\\\hline\end{tabular}\end{center}\begin{tasks}(2)
  \task \rest{1 llanda de A i 2 de B}
  \task $\res{42\text{ €}}$
\end{tasks}

\ex{2} \pau{C. Valenciana, juny 2026}$\min\ C=5x+4y$ \ amb \ $2x+y\geq8$, \ $x+y\geq6$, \ $x+2y\geq10$, \ $x,y\geq0$. Les tres rectes passen per $(2,\,4)$.

\begin{center}\small\begin{tabular}{cc}\hline Vèrtex & $C=5x+4y$\\\hline $(0,\,8)$ & $32$\\$(2,\,4)$ & $26$ \ $\leftarrow$\\$(10,\,0)$ & $50$\\\hline\end{tabular}\end{center}\begin{tasks}(2)
  \task \rest{2 litres de NutriMax i 4 de BioGrow}
  \task $\res{26\text{ €}}$
\end{tasks}


% =====================================================
\seccio{5}{Casos especials: infinites solucions i problemes sense solució}

\ex{2} \begin{center}\small\begin{tabular}{cc}\hline Vèrtex & $B=30x+10y$\\\hline $(0,\,0)$ & $0$\\$(20,\,0)$ & $600$ \ $\leftarrow$\\$(15,\,15)$ & $600$ \ $\leftarrow$\\$(0,\,30)$ & $300$\\\hline\end{tabular}\end{center}El màxim, $\res{600\text{ €}}$, s'assoleix en $(20,\,0)$ i en $(15,\,15)$: $30x+10y=10(3x+y)$ és paral·lela a $3x+y=60$. Hi ha \rest{infinites solucions}: el segment $y=60-3x$, $15\leq x\leq20$. Com que es fabriquen unitats senceres: $(15,\,15)$, $(16,\,12)$, $(17,\,9)$, $(18,\,6)$, $(19,\,3)$ i $(20,\,0)$.

\begin{center}\begin{tikzpicture}[x=0.1105cm,y=0.1105cm,baseline=(current bounding box.center)]\begin{scope}\clip (-3,-3) rectangle (35,35);\foreach \i in {0,10,...,35} \draw[muted!25,very thin] (\i,-3)--(\i,35);\foreach \i in {0,10,...,35} \draw[muted!25,very thin] (-3,\i)--(35,\i);\fill[teal!22] (20.0000,0.0000)--(15.0000,15.0000)--(0.0000,30.0000)--(0.0000,0.0000)--cycle;\end{scope}\draw[-latex,thick] (-3,0)--(37.5,0) node[below left=-1pt,font=\scriptsize]{$x$};\draw[-latex,thick] (0,-3)--(0,37.5) node[below left=-1pt,font=\scriptsize]{$y$};\node[below,font=\tiny,text=muted] at (10,0) {$10$};\node[below,font=\tiny,text=muted] at (20,0) {$20$};\node[below,font=\tiny,text=muted] at (30,0) {$30$};\node[left,font=\tiny,text=muted] at (0,10) {$10$};\node[left,font=\tiny,text=muted] at (0,20) {$20$};\node[left,font=\tiny,text=muted] at (0,30) {$30$};\begin{scope}\clip (-3,-3) rectangle (35,35);\draw[accent,line width=1.1pt] (21.3333,-4.0000)--(8.0000,36.0000);\draw[blue!70!black,line width=1.1pt] (34.0000,-4.0000)--(-6.0000,36.0000);\draw[navy,thick,densely dashed] (21.3333,-4.0000)--(8.0000,36.0000);\draw[accent,line width=3pt] (15.0000,15.0000)--(20.0000,0.0000);\end{scope}\fill[accent] (20.0000,0.0000) circle (2.3pt);\node[font=\tiny,inner sep=1.5pt,fill=white,fill opacity=.75,text opacity=1] at ($(20.0000,0.0000)+(-45.0:0.36cm)$) {$(20,\,0)$};\fill[accent] (15.0000,15.0000) circle (2.3pt);\node[font=\tiny,inner sep=1.5pt,fill=white,fill opacity=.75,text opacity=1] at ($(15.0000,15.0000)+(31.0:0.36cm)$) {$(15,\,15)$};\filldraw[fill=white,draw=navy] (0.0000,30.0000) circle (1.5pt);\node[font=\tiny,inner sep=1.5pt,fill=white,fill opacity=.75,text opacity=1] at ($(0.0000,30.0000)+(115.0:0.36cm)$) {$(0,\,30)$};\filldraw[fill=white,draw=navy] (0.0000,0.0000) circle (1.5pt);\node[font=\tiny,inner sep=1.5pt,fill=white,fill opacity=.75,text opacity=1] at ($(0.0000,0.0000)+(-127.9:0.36cm)$) {$(0,\,0)$};\end{tikzpicture}\end{center}

\ex{2} \begin{tasks}(1)
  \task $f$: $f(0,\,8)=8$, $f(3,\,2)=5$, $f(9,\,0)=9$. Mínim $\res{5}$ en $(3,\,2)$; \rest{no té màxim} (la regió no és acotada i $f$ hi creix indefinidament).
  \task $g$: $g(0,\,8)=8$, $g(3,\,2)=8$, $g(9,\,0)=18$. El mínim, $\res{8}$, s'assoleix en \rest{tot el segment} d'extrems $(0,\,8)$ i $(3,\,2)$ ($g$ és paral·lela a $2x+y=8$). \rest{No té màxim}.
\end{tasks}

\ex{2} Amb $x+y\leq10$ i $x,y\geq0$, el valor més gran de $2x+3y$ és $3\cdot10=30<40$. Les restriccions són incompatibles: la regió factible és \rest{buida} i el problema \rest{no té solució}.

\begin{center}\begin{tikzpicture}[x=0.1739cm,y=0.1739cm,baseline=(current bounding box.center)]\begin{scope}\clip (-1,-1) rectangle (22,15);\foreach \i in {0,5,...,22} \draw[muted!25,very thin] (\i,-1)--(\i,15);\foreach \i in {0,2,...,15} \draw[muted!25,very thin] (-1,\i)--(22,\i);\end{scope}\draw[-latex,thick] (-1,0)--(23.25,0) node[below left=-1pt,font=\scriptsize]{$x$};\draw[-latex,thick] (0,-1)--(0,15.5) node[below left=-1pt,font=\scriptsize]{$y$};\node[below,font=\tiny,text=muted] at (5,0) {$5$};\node[below,font=\tiny,text=muted] at (10,0) {$10$};\node[below,font=\tiny,text=muted] at (15,0) {$15$};\node[below,font=\tiny,text=muted] at (20,0) {$20$};\node[left,font=\tiny,text=muted] at (0,2) {$2$};\node[left,font=\tiny,text=muted] at (0,4) {$4$};\node[left,font=\tiny,text=muted] at (0,6) {$6$};\node[left,font=\tiny,text=muted] at (0,8) {$8$};\node[left,font=\tiny,text=muted] at (0,10) {$10$};\node[left,font=\tiny,text=muted] at (0,12) {$12$};\node[left,font=\tiny,text=muted] at (0,14) {$14$};\begin{scope}\clip (-1,-1) rectangle (22,15);\draw[accent,line width=1.1pt] (12.0000,-2.0000)--(-6.0000,16.0000);\draw[blue!70!black,line width=1.1pt] (-2.0000,14.6667)--(23.0000,-2.0000);\end{scope}\end{tikzpicture}\end{center}

\ex{3} Cal que $f$ siga paral·lela a un costat on hi haja el màxim. Costat $x+y=5$ (de $(5,\,0)$ a $(2,\,3)$): $\res{k=1}$; aleshores $f=x+y$ val 5 en tot el costat i és el màxim. Costat $x-y=-1$ (de $(0,\,1)$ a $(2,\,3)$): $\res{k=-1}$; $f=y-x$ val 1 en tot el costat i és el màxim (en $(0,\,0)$ val 0 i en $(5,\,0)$, $-5$). Per als altres costats (eixos), $f=y$ i $f=x$ no hi tenen el màxim.

\ex{2} \pau{C. Valenciana, juny 1998}\begin{tasks}(1)
  \task $f(0,\,0)=9$, $f(2,\,8)=9$, $f(10,\,3)=-28$. El màxim, $\res{9}$, s'assoleix en \rest{tots els punts del segment} d'extrems $(0,\,0)$ i $(2,\,8)$.
  \task $g(0,\,0)=12$, $g(2,\,8)=28$, $g(10,\,3)=55$: màxim $\res{55}$ en $\res{(10,\,3)}$.
\end{tasks}

\ex{1} \begin{tasks}(1)
  \task \rest{Vertadera}: és el teorema fonamental, i s'assoleixen en vèrtexs.
  \task \rest{Falsa}: en els problemes de dieta la regió no és acotada i hi ha mínim (exercici 17).
  \task \rest{Falsa}: sempre es pot millorar movent-se en la direcció en què creix (o decreix) la funció, fins a arribar a la frontera.
  \task \rest{Vertadera}: també és òptim tot el segment que els uneix.
\end{tasks}


\end{document}
